\RequirePackage{plautopatch}
\documentclass[14pt,aspectratio=169,xcolor=dvipsnames,table,onlytextwidth,dvipdfmx]{beamer}
\usepackage[utf8]{inputenc}
\usepackage{bxdpx-beamer} % dvipdfmxなので必要

\usetheme{Boadilla}
%Beamerフォント設定
\usefonttheme{professionalfonts}
\usepackage[T1]{fontenc}
\usepackage[deluxe,uplatex]{otf} % 日本語多ウェイト化
\usepackage{mlmodern}  % 太いComputer Modern
% MLmodernのバグを修正: cf. https://tex.stackexchange.com/questions/646333/size-of-integral-symbol-in-section-header-with-mlmodern
\DeclareFontFamily{OMX}{mlmex}{}
\DeclareFontShape{OMX}{mlmex}{m}{n}{%
   <->mlmex10%
   }{}
\usepackage{newtxtext} % 数式以外をTXフォントで上書き
% \usepackage[noalphabet,unicode]{pxchfon}
% \setminchofont{KozMinPro-Regular.otf}% 游明朝Regular
% \setboldminchofont{KozMinPro-Bold.otf}% 游明朝Demibold
% \setgothicfont[0]{BIZ-UDGothicR.ttc}% BIZ UDゴシックR
% \setboldgothicfont[0]{BIZ-UDGothicB.ttc}% BIZ UDゴシックB
\renewcommand{\familydefault}{\sfdefault}  % 英文をサンセリフ体に
\renewcommand{\kanjifamilydefault}{\gtdefault}  % 日本語をゴシック体に
\usefonttheme{structurebold} % タイトル部を太字
\setbeamerfont{alerted text}{series=\bfseries} % Alertを太字
\setbeamerfont{section in toc}{series=\mdseries} % 目次は太字にしない
\setbeamerfont{frametitle}{size=\Large} % フレームタイトル文字サイズ
\setbeamerfont{title}{size=\LARGE} % タイトル文字サイズ
\setbeamerfont{date}{size=\small}  % 日付文字サイズ
\usepackage{bxcoloremoji}

% Babel (日本語の場合のみ・英語の場合は不要)
\uselanguage{japanese}
\languagepath{japanese}
\deftranslation[to=japanese]{Theorem}{定理}
\deftranslation[to=japanese]{Lemma}{補題}
\deftranslation[to=japanese]{Example}{例}
\deftranslation[to=japanese]{Examples}{例}
\deftranslation[to=japanese]{Definition}{定義}
\deftranslation[to=japanese]{Definitions}{定義}
\deftranslation[to=japanese]{Problem}{問題}
\deftranslation[to=japanese]{Solution}{解}
\deftranslation[to=japanese]{Fact}{事実}
\deftranslation[to=japanese]{Proof}{証明}
\deftranslation[to=japanese]{Corollary}{系}
\setbeamertemplate{theorems}[normal font] %日本語の場合，定理内を斜体にしない
% overlay-aware proof environment
\renewenvironment<>{proof}{%
\par
\begin{actionenv}#1%
\structure{(証明)}}
{\qed\end{actionenv}}

%Beamer色設定
\definecolor{UniBlue}{RGB}{0,150,200}
\definecolor{UniGreen}{RGB}{3,175,122}
\definecolor{UniPink}{RGB}{255,128,255}
\definecolor{AlertOrange}{RGB}{255,76,0}
\definecolor{AlmostBlack}{RGB}{38,38,38}
\setbeamercolor{normal text}{fg=AlmostBlack}  % 本文カラー
\setbeamercolor{structure}{fg=UniBlue} % 見出しカラー
\setbeamercolor{block title}{fg=UniBlue!50!black} % ブロック部分タイトルカラー
\setbeamercolor{alerted text}{fg=AlertOrange} % \alert 文字カラー
\setbeamercolor{bibliography entry author}{fg=AlmostBlack} % Bib 文字カラー
\setbeamercolor{bibliography entry note}{fg=AlmostBlack} % Bib 文字カラー
\mode<beamer>{
    \definecolor{BackGroundGray}{RGB}{254,254,254}
    \setbeamercolor{background canvas}{bg=BackGroundGray} % スライドモードのみ背景をわずかにグレーにする
}

%フラットデザイン化
\setbeamertemplate{blocks}[rounded] % Blockの影を消す
\useinnertheme{circles} % 箇条書きをシンプルに
\setbeamertemplate{navigation symbols}{} % ナビゲーションシンボルを消す
\setbeamertemplate{footline}[frame number] % フッターはスライド番号のみ

%タイトルページ
\setbeamertemplate{title page}{%
    \vspace{2.5em}
    {\usebeamerfont{title} \usebeamercolor[fg]{title} \inserttitle \par}
    {\usebeamerfont{subtitle}\usebeamercolor[fg]{subtitle}\insertsubtitle \par}

    \vspace{1.5em}
    \begin{columns}[]
        \begin{column}{.35\linewidth}
        \usebeamerfont{titlegraphic}\inserttitlegraphic\par
        \end{column}
        \begin{column}{.6\linewidth}\setlength\topsep{0pt}
            \begin{flushright}
        \usebeamerfont{author}\insertauthor\par
        \usebeamerfont{institute}\insertinstitute \par
        \vspace{3em}
        \usebeamerfont{date}\insertdate\par
            \end{flushright}
        \end{column}
    \end{columns}
}

%biblatex
% \usepackage[style=ext-authoryear,citestyle=authoryear,natbib=true,giveninits=true,maxbibnames=99,maxcitenames=10,url=false,isbn=false,doi=false,articlein=false]{biblatex}
% \DeclareNameAlias{author}{first-last}
% \addbibresource{shrunk.bib}
% \newcommand{\mkbibbracketscol}[1]{\textcolor{UniBlue}{\scriptsize \mkbibbrackets{#1}}}
% \DeclareCiteCommand{\cite}[\mkbibbracketscol]
%   {\usebibmacro{prenote}}
%   {\usebibmacro{citeindex}%
%    \usebibmacro{cite}}
%   {\multicitedelim}
%   {\usebibmacro{postnote}}
% % \renewcommand{\cite}[1]{\textcolor{UniBlue}{\scriptsize\citep{#1}}}
% \renewcommand*{\finalnamedelim}{\addcomma\addspace} % remove "and" before the last author
% \DeclareDelimFormat{nameyeardelim}{\addspace}
% \setbeamertemplate{bibliography item}{}

% Algorithm系
\usepackage{algorithm}
\usepackage[noend]{algorithmic}
\algsetup{linenosize=\color{fg!50}\footnotesize}
\renewcommand\algorithmicdo{:}
\renewcommand\algorithmicthen{:}
\renewcommand\algorithmicrequire{\textbf{Input:}}
\renewcommand\algorithmicensure{\textbf{Output:}}
\renewcommand\algorithmiccomment[1]{\hfill\textcolor{fg!50}{\footnotesize\textit{// #1}}}

% TikZ
\usepackage{tikz}
\usetikzlibrary{positioning,shapes,arrows,quotes,shapes.callouts,calc,graphs,graphs.standard,fit,backgrounds,perspective,matrix}
\tikzset{point/.style={circle, fill=AlertOrange, inner sep=1pt, text width=1pt}}
\tikzset{bpoint/.style={circle, fill=black, inner sep=1pt, text width=1pt}}
\tikzset{alt/.code args={<#1>#2#3}{\alt<#1>{\pgfkeysalso{#2}}{\pgfkeysalso{#3}}}}
\tikzset{>=latex}
\tikzset{mynodes/.style={circle,white,fill=UniBlue!90,text width=.8em,inner sep=0pt,text centered,font=\footnotesize}}
\tikzset{myedges/.style={thick}}
\tikzset{mypath/.style={ultra thick, draw=#1}, mypath/.default=AlertOrange}
\tikzset{mysubset/.style={draw,#1,dotted,thick,fill=#1!10}, mysubset/.default=UniBlue}
\tikzset{mypath2/.style={ultra thick, dashed, draw=UniGreen}}
\tikzset{graphs/mygraph/.style={nodes={mynodes}, edges={myedges}, empty nodes}}
\tikzset{nomargin/.style={inner sep=0pt, outer sep=0pt}}
%% graph macros
\tikzgraphsset{declare={mybipartite}{%
subgraph I_nm [V={1,2,3,4,5,6}, W={a,b,c,d,e,f}];
1 -- {a,b};
2 -- {b,c,d,e};
3 -- {d,f};
4 -- {e,f};
{5,6} -- f;
}}
\tikzgraphsset{declare={mybipartitedi}{%
subgraph I_nm [V={1,2,3,4,5,6}, W={a,b,c,d,e,f}];
1 -> {a,b};
2 -> {b,c,d,e};
3 -> {d,f};
4 -> {e,f};
{5,6} -> f;
}}
\tikzgraphsset{declare={mybipartite2}{%
subgraph I_nm [V={1,2,3}, W={a,b,c}];
1 -- {a,b}; 2 -- {a,b,c}; 3 -- {a,c};
}} 
% Petersen graph
\tikzgraphsset{declare={Petersen}{%
    subgraph C_n [n=5,name=A, radius=1.5cm]; 
    subgraph I_n [n=5,name=B, radius=.8cm];
    \foreach \i [evaluate={\j=int(mod(\i+2,5)+1);}] in {1,...,5}{
        A \i -- B \i;
        B \i -- B \j;
}}}
% Example graph for Menger's theorem
\tikzgraphsset{declare={menger}{%
    s[at={(-1,0)}, as={$s$}];
    (s) -> 1[at={(0,1)}] -> 2[at={(1,0)}];
    (s) -> 3[at={(0,-1)}] -> 2;
    4[at={(2,1)}] <- 5[at={(2,-1)}];
    {5,6[at={(3,1)}], 7[at={(3,-1)}]} -> t[at={(4,0)}, as={$t$}];
    (s) -> (2);
    (2) -> {(4), (5)};
    {(4),(5)} -> {(6),(7)};
    (6) -> (7);
    1 -> 4;
}} 
\tikzgraphsset{declare={undirected menger}{%
    s[at={(-1,0)}, as={$s$}];
    (s) -- 1[at={(0,1)}] -- 2[at={(1,0)}];
    (s) -- 3[at={(0,-1)}] -- 2;
    4[at={(2,1)}] -- 5[at={(2,-1)}];
    {5,6[at={(3,1)}], 7[at={(3,-1)}]} -- t[at={(4,0)}, as={$t$}];
    (s) -- (2);
    (2) -- {(4), (5)};
    {(4),(5)} -- {(6),(7)};
    (6) -- (7);
    1 -- 4;
}} 

% Appendix
\usepackage{appendixnumberbeamer}

% macros
\newcommand{\R}{\mathbb{R}}
\newcommand{\Z}{\mathbb{Z}}
\newcommand{\be}{\mathbf{e}}
\newcommand{\zeros}{\mathbf{0}}
\newcommand{\ones}{\mathbf{1}}
\newcommand{\bfzero}{\zeros}
\newcommand{\bfone}{\ones}

% \newcommand{\bp}{\mathbf{p}}
\newcommand{\eps}{\varepsilon}
\newcommand{\defiff}{\overset{\text{def}}{\iff}}
\DeclareMathOperator{\poly}{poly}
\DeclareMathOperator{\conv}{conv}
\DeclareMathOperator{\In}{In}
\DeclareMathOperator{\Out}{Out}
\DeclareMathOperator{\val}{val}
\newcommand{\lo}{{\mathrm{lo}}}
\newcommand{\up}{{\mathrm{up}}}
\newcommand{\caI}{\mathcal{I}}
\newcommand{\caB}{\mathcal{B}}
\newcommand{\caC}{\mathcal{C}}
\newcommand{\bM}{\mathbf{M}}
\newcommand{\IO}{\mathrm{IO}}

\usepackage{mathtools}
\DeclarePairedDelimiter{\abs}{\lvert}{\rvert}
\DeclarePairedDelimiter{\norm}{\lVert}{\rVert}
\DeclarePairedDelimiter{\inprod}{\langle}{\rangle}
\usepackage[no-test-for-array]{nicematrix}
\usepackage{diagbox}

\newcommand{\highlightcap}[3][red]{\tikz[baseline=(x.base)]{\node[rectangle,rounded corners,fill=#1!10](x){#2} node[below=0pt of x, color=#1]{#3};}}
\newcommand{\highlight}[2][red]{\tikz[baseline=(x.base)]{\node[rectangle,rounded corners,fill=#1!10](x){#2};}}
\newcommand{\mycite}[2][\footnotesize]{\textcolor{UniBlue}{#1 [#2]}}

%section
% \AtBeginSection[]{
%     \frame[c]{
%         \centering
%         \begin{tikzpicture}
%             \node[white, circle, fill=UniBlue, text centered, text width=.3\linewidth](bullet) {
%                 {\Huge \bf \insertsectionnumber} \\[4em]
%     };
%     \node[text width=.65\linewidth, right=1em of bullet]{%
%         \usebeamerfont{title}\usebeamercolor[fg]{title}\insertsection\par %
%     };
%         \end{tikzpicture}
%     } %目次スライド
% }

%grid
% \usepackage[colorgrid,gridunit=pt,texcoord]{eso-pic}
\newcommand{\postit}[1]{\tikz[overlay, remember picture]{\node[draw,anchor=north east, align=center, rectangle, fill=yellow!10, xshift=-1pt, yshift=-1pt] at (current page.north east) {#1};}}

\usepackage{tcolorbox}
\newtcolorbox{primal}{colback=red!10!white, colframe=red, title=主問題 (P), left=2pt,right=2pt,top=0pt,bottom=0pt}
\newtcolorbox{dual}{colback=blue!10!white, colframe=blue, title=双対問題 (D), left=2pt,right=2pt,top=0pt,bottom=0pt}

% visible on
\usetikzlibrary{overlay-beamer-styles}
\usepackage[normalem]{ulem}
\usepackage{pxrubrica}
\usepackage{booktabs}

\newcommand{\Aphiup}{\textcolor{UniBlue}{A_\varphi^\up}}
\newcommand{\Aphilow}{\textcolor{AlertOrange}{A_\varphi^\lo}}
\newcommand{\bv}{\boldsymbol{v}}
\newcommand{\bx}{\boldsymbol{x}}
\newcommand{\by}{\boldsymbol{y}}
\newcommand{\bz}{\boldsymbol{z}}
\newcommand{\bw}{\boldsymbol{w}}
\newcommand{\bc}{{\boldsymbol c}}
\DeclareMathOperator{\supp}{supp}
\DeclareMathOperator{\argmax}{argmax}
\DeclareMathOperator{\argmin}{argmin}
\DeclareMathOperator{\E}{\mathbb{E}}
\csname endofdump\endcsname
\title{組合せ最適化特論}
\subtitle{第6回（最小費用流）}
\author{担当: \textbf{相馬 輔}}
\date{2025/12/23}

\begin{document}
\maketitle

\begin{frame}{}
    \centering\small
    \tikzset{block/.style={rectangle, draw=gray, fill=gray!10, thick, minimum width=3em, minimum height=2em, align=center}}
    \begin{tikzpicture}
        \node[block,] (BM) {重みなし\\二部マッチング};
        \node[block,above=1em of BM.north west, xshift=-2em] (WBM) {重みつき\\二部マッチング};
        \node[block,above=5em of BM.north east, xshift=1em] (NF) {最大流};
        \node[block,above=9em of BM] (MCF) {最小費用流};
        % \node[block,above=1em of NF.north east, xshift=1em] (MC) {最小カット};
        \node[block,right=11em of BM] (MTR) {マトロイド};
        \node[block,above=7em of MTR] (MI) {マトロイド交差};
        \node[block,above=3em of MCF, xshift=7em] (SFM) {劣モジュラ最小化};
        \graph[mygraph]{
            (BM) -> {(WBM), (NF)} -> (MCF);
            {(MTR), (BM)} -> (MI);
            % (NF) -> (MC);
            {(MCF), (MI)} -> (SFM);
        };
        \begin{scope}[on background layer]
            \node[mysubset, fit=(BM) (WBM), "二部マッチング" {below, UniBlue}]{};
            \node[mysubset, AlertOrange, fill=AlertOrange!10, fit=(NF) (MCF), "ネットワークフロー" {above left, AlertOrange}]{};
            \node[mysubset, UniGreen, fill=UniGreen!10, fit=(MTR) (MI), "マトロイド" {below, UniGreen}]{};
        \end{scope}
    \end{tikzpicture}
\end{frame}

\begin{frame}[allowframebreaks]
    \frametitle{各回の内容（予定）} 
    \small
    \begin{enumerate}
        \item (4/23) イントロ+多面体的組合せ論（線形計画法の復習，整数多面体，完全単模行列）
        \item[] \textcolor{gray}{(4/30) 休み}
        \item (5/7) 二部マッチング①（Konig-Egervaryの定理，増加道アルゴリズム，ハンガリー法）
        \item (5/14) 二部マッチング②（最短路問題の復習，逐次最短路法と主双対法，最適性基準からの見方）
        \item (5/21) 最大流① （定式化，最大流最小カット定理，応用例）
        \item (5/28) 最大流②（残余ネットワーク，Ford--Fulkerson法，Edmonds--Karpのアルゴリズム）+ 最小費用流①（定式化，輸送問題，最大流との関係）
        \item[] \textcolor{gray}{(6/4) 休み}
        % \item[] \sout{(12/11) 最小カット（永持茨木のアルゴリズム，Kargerのアルゴリズム）}
        \item (6/11) 最小費用流②（逐次最短路法，容量スケーリング法）
        \item[] \textcolor{gray}{(6/18) 休み}
        \item[] \textcolor{gray}{(6/25) 休み}
        \item (7/2) マトロイド（定義と公理系，貪欲法，マトロイド多面体） 
        \item (7/9) 独立マッチング・マトロイド交差 （定義，応用例，Edmondsの最大最小定理，増加道アルゴリズム）
        \item (7/16) 劣モジュラ関数①（諸例，劣モジュラ基多面体，Lov\'asz拡張，劣モジュラ最小化）
        \item[] \textcolor{gray}{(7/23) 休み}
        \item (7/30) 劣モジュラ関数②（劣モジュラ最大化，近似アルゴリズム，貪欲法）
        \item (どこか) 精選トピック①
        \item (どこか) 精選トピック②
        \item (どこか) 精選トピック③
    \end{enumerate}

    \bigskip
    \begin{itemize}
        \item レポート課題を6月と期末に出す予定
        \item 精選トピックの日時は今後調整 
    \end{itemize}
\end{frame}

\begin{frame}<1>[label=outline]{目次}
\begin{tikzpicture}
    \tikzstyle{block} = [rectangle, text width=\textwidth, rounded corners, minimum height=4em]
    \tikzstyle{line} = [draw, -latex]

    \node[block,alt={<1,3->{fill=cyan!10}{fill=red!10}}] (combopt) {
        \begin{columns}
            \begin{column}{0.5\linewidth}
        \textbf{1. 最小費用流問題}
            \end{column}
            \begin{column}{0.4\linewidth}
                \small
                \structure{・} 定式化 \\
                \structure{・} Hitchcock型輸送問題 \\
                \structure{・} Galeの定理 \\
                \structure{・} 負閉路最適性条件 \\
            \end{column}
        \end{columns}
    };
    \node[block,alt={<1-2,4->{fill=cyan!10}{fill=red!10}}, below=.5em of combopt] (LP) {
        \begin{columns}
            \begin{column}{0.5\linewidth}
        \textbf{2. アルゴリズム}
            \end{column}
            \begin{column}{0.4\linewidth}
                \small
                \structure{・} 逐次最短路法，主双対法 \\
                \structure{・} 容量スケーリング \\
                \structure{・} まとめ
            \end{column}
        \end{columns}
    };
\end{tikzpicture}%
\end{frame}


%%%%%%%%%%%%%%%%%%%%%%%%%%%%%%%%%%%%%%%%%%%%%
\section{最小費用流問題}
\againframe<2|handout:0>{outline}
%%%%%%%%%%%%%%%%%%%%%%%%%%%%%%%%%%%%%%%%%%%%%
\begin{frame}{境界}
    $G = (V, A)$: 有向グラフ 

    \begin{definition}
        $\varphi: A \to \R$に対し，その\alert{境界} $\partial \varphi: V \to \R$を
        \[
            (\partial\varphi)(i) := \sum_{a \in \Out(i)} \varphi(a) - \sum_{a \in \In(i)} \varphi(a)  \qquad (i \in V)
        \]
        と定義する．ここで，$\In(i)$, $\Out(i)$はそれぞれ頂点$i$に入る枝，出る枝の集合．

        {\hfill\footnotesize ※$\delta^+(i)$, $\delta^-(i)$と表記する本もある．}
    \end{definition}
\end{frame}

\begin{frame}{$b$流}
    $G = (V, A)$: 有向グラフ, $u: A \to \R_+$: 枝容量関数, \\
    \alert{$b: V \to \R$ with $b(V) = 0$}

    \begin{definition}[$b$流]
        $\varphi: A \to \R$が\alert{（実行可能）$b$流} $\defiff$
        \begin{itemize}
            \item $0 \leq \varphi(a) \leq u(a)$ \qquad ($a \in A$) \hfill \textbf{（容量制約）} 
            \item $(\partial\varphi)(i) = b(i)$ \qquad ($i \in V$) \hfill \textbf{（境界条件）}
        \end{itemize}
    \end{definition}

    \structure{解釈}
    \begin{itemize}
        \item $b(i) > 0$: 供給点，湧き出し口
        \item $b(i) = 0$: 流量保存する頂点
        \item $b(i) < 0$: 需要点，吸い込み口
    \end{itemize}

    特に，$b \equiv 0$のとき，$\varphi$は\alert{循環流 (circulation)}と呼ばれる．
\end{frame}

\begin{frame}{最小費用流問題}

    \begin{itemize}
    \item $G = (V, A)$: 有向グラフ
    \item $u: A \to \R_+$: 枝容量関数 \\ 
    \item $b: V \to \R$ with $b(V) = 0$ \\ 
    \item \alert{$c: A \to \R$: 枝費用関数}
    \end{itemize}

    \bigskip
    \begin{block}{最小費用流問題 (Minimum Cost Flow Problem)}
        $b$流$\varphi$で費用$c(\varphi) := \sum_{a \in A} c(a) \varphi(a) $が最小のものを求めよ． \\
        （もしくは，$b$流が存在しないことを示せ）
    \end{block}
\end{frame}

\begin{frame}<1>[label=mcf-LP]{最小費用流: 線形計画問題として}
    最小費用流の方がLPとして美しい形になる．

    \footnotesize
    \begin{primal}
        \setlength{\abovedisplayskip}{0pt}
        \begin{alignat*}{3}
            &\text{min}   &\quad& c^\top\varphi &\quad& \\
            &\text{s.t.}  &\quad& \sum_{a \in \Out(i)} \varphi(a)  - \sum_{a \in \In(i)} \varphi(a) = b(i) &\quad& (i \in V) \\
            &             &\quad& 0 \leq \varphi(a) \leq u(a) &\quad& (a \in A) 
        \end{alignat*}
    \end{primal}
    \onslide<2>{（後の都合のため）$x$の符号を反転すると．．．}
    \begin{dual}
        \setlength{\abovedisplayskip}{0pt}
        \alt<2>{
        \begin{alignat*}{3}
            &\max_{x \in \R^V, y \in \R_+^A}  &\quad& -\sum_{i \in V} b(i) x(i) - \sum_{a \in A} u(a) y(a) &\quad& \\
            &\text{s.t.} &\quad& -x(i) + x(j) - y(a) \leq c(a)  &\quad& (a=ij \in A) \\
            &            &\quad& y(a) \geq 0  &\quad& (a \in A)
        \end{alignat*}
        }{
        \begin{alignat*}{3}
            &\max_{x \in \R^V, y \in \R_+^A}  &\quad& \sum_{i \in V} b(i) x(i) - \sum_{a \in A} u(a) y(a) &\quad& \\
            &\text{s.t.} &\quad& x(i) - x(j) - y(a) \leq c(a)  &\quad& (a=ij \in A) \\
            &            &\quad& y(a) \geq 0  &\quad& (a \in A)
        \end{alignat*}
        }
    \end{dual}
\end{frame}


\tikzgraphsset{declare={flowexample}{%
    s[at={(0,0)}, as={$s$}];
    a[at={(2, 1)}];
    b[at={(2,-1)}];
    c[at={(4, 1)}];
    d[at={(4,-1)}];
    t[at={(6, 0)}, as={$t$}];
}}

\begin{frame}{最大流 $\subset$ 最小費用流}
    $s$--$t$最大流は最小費用流問題の特殊ケース．

    \bigskip
    \small
    \begin{columns}[T]
    \begin{column}{.6\textwidth}
    $G$に枝$a^* = ts$を追加した有向グラフを$G'$とし，その上の循環流を考える．

    \bigskip 
    容量
    \[
        u'(a) := \begin{cases}
            +\infty & (a = a^*) \\
            u(a) & (a \in A) 
        \end{cases} 
    \]
    枝費用
    \[
        c(a) := \begin{cases}
            -1 & (a = a^*) \\
            0 & (\text{otherwise}) 
        \end{cases}
    \]
    とすればよい．
    \end{column}
    \begin{column}{.4\textwidth}
        \centering
    \begin{tikzpicture}[xscale=.7]
        \graph[no placement, mygraph]{flowexample};
        \foreach \i/\j/\f/\u in {
            s/a/5/5,
            s/b/2/3,
            a/b/2/2,
            a/c/3/3,
            d/a/0/2,
            b/d/4/4,
            c/t/4/5,
            d/c/1/1,
            d/t/3/3
        } {
            \path[draw,myedges,->] (\i) -- (\j);
        }
        \coordinate[below=4em of t] (tbelow);
        \draw[AlertOrange,mypath,->] (t) -- (tbelow) -- node[below]{$a^*$} (s |- tbelow) -- (s);
    \end{tikzpicture}
    \end{column}
    \end{columns}
\end{frame}

\begin{frame}{最短路問題 $\subset$ 最小費用流}
    最短路問題は最小費用流問題の特殊ケース．

    \bigskip
    \small
    \begin{itemize}
        \item 容量: $u \equiv +\infty$ 
        \item 境界:
    $b(i) = 
        \begin{cases}
            +1 & (i = s) \\
            -1 & (i = t) \\
            0  & (\text{otherwise})
        \end{cases}$ 
        \item 費用 $c = \ell$ (最短路問題の枝長)
    \end{itemize}
\end{frame}

\begin{frame}{Hitchcock型輸送問題}
    \small
    $G = (V^+, V^-; E)$: 二部グラフ \\
    $b^+: V^+ \to \R_+$: 供給関数, $b^-: V^- \to \R_+$: 需要関数 with $b^+(V^+) = b^-(V^-)$ \\
    $c: E \to \R$: 枝費用関数

    \begin{columns}[T]
    \begin{column}{.5\textwidth}
    \begin{block}{Hitchcock型輸送問題}
        \setlength{\abovedisplayskip}{0pt}
        \begin{alignat*}{3}
            &\text{min}   &\quad& \sum_{(i,j) \in E} c(i,j) p(i,j)  & \\
            &\text{s.t.} &\quad& \sum_{j \in \Gamma(i)} p(i, j) = b^+(i) &\quad& (i \in V^+) \\
            &                  &\quad& \sum_{i \in \Gamma(j)} p(i, j) = b^-(j) &\quad& (j \in V^-) \\
            &                  &\quad& p(i,j) \geq 0 &\quad& ((i,j) \in E)
        \end{alignat*}
    \end{block}
    \end{column}
    \begin{column}{.5\textwidth}
        \centering
        \begin{tikzpicture}[>=latex]
        \foreach \i in {1,2,3} {
            \node[inner sep=0pt](v\i) at ($(0, -\i*2)$) {\coloremojiucs[scale=2]{:factory:}};
            \node[inner sep=0pt](w\i) at ($(3, -\i*2)$) {\coloremojiucs[scale=2]{:convenience_store:}};
        }
        \graph[mygraph]{
            {(v1), (v2), (v3)} --[complete bipartite] {(w1), (w2), (w3)}
        };
        \path (v1)--(w1) node[midway,above,align=center]{\footnotesize 輸送 $p(i,j)$ \\ \reflectbox{\coloremoji{🚛}}};
        \node[above=5pt of v1] {$V^+$};
        \node[above=5pt of w1] {$V^-$};
        \node[left =3pt of v1, font=\footnotesize, align=center] {供給 \\ $b^+(i)$};
        \node[right=3pt of w1, font=\footnotesize, align=center] {需要 \\ $b^-(j)$};
        \end{tikzpicture}
    \end{column}
    \end{columns}
\end{frame}

\begin{frame}{輸送問題と最小費用流}
    \small
    \begin{lemma}
        $n$頂点, $m$枝の最小費用流問題は，$n+m$頂点, $2m$枝のHitchcock型輸送問題に帰着できる．
    \end{lemma}

    \begin{columns}[T]
    \begin{column}{.65\textwidth}
    \structure{(証明)}
    $G' = (V^+, V^-; E)$, $b^+$, $b^-$, $c$を次のように構成する:
    \begin{itemize}
        \setlength{\itemsep}{0em}
        \item $V^+ := A$, $V^- := V$
        \item $E = \bigcup_{a=ij \in A} \{ (a, i), (a, j)\}$
        \item $b^+(a) := u(a)$ \quad ($a \in V^+$) \\
              $b^-(i) := u(\Out(i)) - b(i)$ \quad ($i \in V^-$)
        \item $c(a, i) := 0$, $c(a, j) := c(a)$ \quad ($a=ij \in A$)
    \end{itemize}
    \end{column}
    \begin{column}{.35\textwidth}
        \centering
        \begin{tikzpicture}
            \graph[mygraph, grow right=1.7cm]{
                subgraph I_nm[V={a1, a2}, W={1, 2, 3}];
                a1 -- {1, 2};
                a2 -- {2, 3};
            };
            \node[above=1pt of a1, font=\footnotesize, align=center] {$A$};
            \node[above=1pt of 1, font=\footnotesize, align=center] {$V$};
            \node[left=1pt of a1, font=\footnotesize, align=center] {$u(a)$};
            \node[right=1pt of 2, font=\footnotesize, align=center] {$u(\Out(i)) - b(i)$};
        \end{tikzpicture}
    \end{column}
    \end{columns}
    \bigskip
    $b^+(V^+) = u(A) = \sum_{i \in V} (u(\Out(i)) - b(i)) = b^-(V^-)$より，これはHitchcock型輸送問題を定める．
\end{frame}

\begin{frame}{輸送問題と最小費用流}
    \small
    \begin{lemma}
        $n$頂点, $m$枝の最小費用流問題は，$n+m$頂点, $2m$枝のHitchcock型輸送問題に帰着できる．
    \end{lemma}

    \begin{columns}[T]
    \begin{column}{.65\textwidth}
    \structure{(証明)}
    $(G, u)$上の$b$流$\varphi$に対し，
    \begin{align*}
        p(a, i) &:= u(a) - \varphi(a) \\
        p(a, j) &:= \varphi(a) 
    \end{align*}
    ($a=ij \in A$)と定めると，$p$は輸送問題の実行可能解であり，$c'(p) = c(\varphi)$.
    \end{column}
    \begin{column}{.35\textwidth}
        \centering
        \begin{tikzpicture}
            \graph[mygraph, grow right=1.7cm]{
                subgraph I_nm[V={a1, a2}, W={1, 2, 3}];
                a1 -- {1, 2};
                a2 -- {2, 3};
                {[edges={UniBlue, thick}]
                    a1 --["{\footnotesize $u(a) - \varphi(a)$}", above] 1; 
                    a1 --["{\footnotesize $\varphi(a)$}", {below, sloped}] 2; 
                }
            };
            \node[above=6pt of a1, font=\footnotesize, align=center] {$A$};
            \node[above=6pt of 1, font=\footnotesize, align=center] {$V$};
            \node[left=1pt of a1, font=\footnotesize, align=center] {$u(a)$};
            \node[right=1pt of 2, font=\footnotesize, align=center] {$u(\Out(i)) - b(i)$};
        \end{tikzpicture}
    \end{column}
    \end{columns}
    \bigskip
    逆に，輸送問題の実行可能解$p$に対し，
    \begin{align*}
        \varphi(a) := p(a, j) \quad (a=ij \in A)
    \end{align*}
    と定めると，$\varphi$は$(G, u)$上の$b$流であり，$c(\varphi) = c'(p)$. \qed
\end{frame}

\begin{frame}{実行可能性: Galeの定理}
    \small
    \begin{theorem}
        （もし存在するならば）$b$流は1回の最大流問題で求められる．
    \end{theorem}
    \begin{columns}
    \begin{column}{.7\textwidth}
    \structure{(証明)}
    右図のグラフ$G' = (V', A')$と枝容量を考える:
    \begin{align*}
        u'(a) &= 
        \begin{cases}
            u(a)  & (a \in A) \\
            b(i)  & (a=si, b(i) > 0) \\
            -b(i) & (a=it, b(i) < 0) \\
        \end{cases}
    \end{align*}
    $(G', u')$の流量$b(V^+)$の$s$--$t$流 $\longleftrightarrow$ $(G, u)$の$b$流 \qed
    \end{column}
    \begin{column}{.3\textwidth}
    
    \end{column}
    \end{columns}

    \pause\bigskip
    \begin{corollary}[Gale, 1957]
        $b$流が存在 $\iff$ $b(X) \leq u(\Out_G(X))$ \quad ($X \subseteq V$)
    \end{corollary}
    \structure{(証明)} ($G'$の$s$--$t$カット容量) $\geq b(V^+)$ を書き換える． \qed
\end{frame}

\begin{frame}{残余ネットワーク}
    \small
    $b$流に関しても残余ネットワークの定義は同じ．

    \begin{definition}[残余ネットワーク]
        \setlength{\abovedisplayskip}{2pt}
        \setlength{\belowdisplayskip}{2pt}
        \begin{itemize}
            \item 
    $b$流$\varphi$に対し，\alert{残余ネットワーク} $G_\varphi = (V, A_\varphi)$を次のように定める: 
    \begin{align*}
        A_\varphi &:= \Aphiup \cup \Aphilow \\
        \Aphiup &:= \{ a \in A : \varphi(a) < u(a) \} \quad\text{\color{gray}\footnotesize \dots\dots\dots$\varphi$を増やせる枝} \\
        \Aphilow &:= \{ a^{-1} : a \in A,\; \varphi(a) > 0  \} \quad \text{\color{gray}\footnotesize \dots\dots\dots$\varphi$を減らせる枝}
    \end{align*}
    ここで$a^{-1}$は$a$の逆向き枝．

    \item 残余ネットワークの枝容量$c_\varphi: A_\varphi \to \R_+$を
    次のように定める:
    \begin{align*}
        c_\varphi(a) := \begin{cases}
            u(a) - \varphi(a) & (a \in \Aphiup) \\
            \varphi(a^{-1})        & (a \in \Aphilow)
        \end{cases}
    \quad\text{\textcolor{gray}{\footnotesize \dots\dots\dots$\varphi$を増減できる余裕量}}
    \end{align*}
        \end{itemize}
    \end{definition}
\end{frame}

\begin{frame}{残余ネットワークを用いた流の更新}
    \small
    $\varphi$: $b$流 
    
    \bigskip
    残余ネットワーク$G_\varphi$上のパス or 閉路$Q$に沿って流を更新する:
    \[
        \varphi^+(a) := \begin{cases}
            \varphi(a) \textcolor{UniBlue}{{}+\eps} & (a \in Q \cap \Aphiup) \\
            \varphi(a) \textcolor{AlertOrange}{{}-\eps} & (a^{-1} \in Q \cap \Aphilow) \\
            \varphi(a)        & (\text{otherwise})
        \end{cases}
    \]
    ここで，$\eps \leq \min_{a \in Q} c_\varphi(a)$.

    \pause
    \bigskip
    \structure{観察}
        \begin{itemize}
            \item $Q$が$s$--$t$パスならば，$\varphi^+$は$b' = b + \eps(\be_s - \be_t)$に関する$b'$流．
            \item $Q$が閉路ならば，$\varphi^+$は$b$流．
        \end{itemize}
\end{frame}

\begin{frame}{残余ネットワーク}
    \small
    さらに残余ネットワークの枝長$\ell_\varphi: A_\varphi \to \R$を次のように定める:
        \[
            \ell_\varphi(a) := \begin{cases}
                \phantom{-}c(a) & (a \in \Aphiup) \\
                -c(a^{-1}) & (a \in \Aphilow)
            \end{cases}
        \]
    \structure{観察} $Q$に沿った流の更新により，費用は次のように変化する:
    \begin{block}{}
    \setlength{\abovedisplayskip}{0em}
    \begin{align*}
        c(\varphi^+) = c(\varphi) + \eps \ell_\varphi(Q)
    \end{align*}
    \end{block}

    \vfill
    \pause\vfill
    \structure{(復習)}
    $p: V \to \R$が枝長$\ell: A \to \R$に関する\alert{ポテンシャル} \\ 
    $\defiff$ $p(i) - p(j) \leq \ell(a)$ \quad ($a=ij \in A$)
\end{frame}

\begin{frame}{負閉路最適性条件}
    \begin{theorem}[Klein (1967), Ford--Fulkerson (1962)]
        $b$流$\varphi$に対し，次の3条件は同値:
        \begin{enumerate}
            \item $\varphi$は最小費用流
            \item 残余ネットワーク$G_\varphi$に（$\ell_\varphi$に関する）負閉路が存在しない
            \item 残余ネットワーク$G_\varphi$に（$\ell_\varphi$に関する）ポテンシャルが存在
        \end{enumerate}
    \end{theorem}

    最小重み完全マッチングの最適性条件（第3回）の一般化．

    \pause
    \bigskip
    \small
    \structure{(証明)}
    \textbf{①$\implies$②:}
    負閉路$C$に沿って$\varphi$を更新すると，$\varphi^+$は$b$流で，$c(\varphi^+) = c(\varphi) + \eps \ell_\varphi(C) < c(\varphi)$となり矛盾．
    
    \pause
    \bigskip
    \textbf{②$\iff$③:} 最短路問題（第3回）でやった

    \pause
    \bigskip
    \textbf{③$\implies$①:} LPの相補性条件から従う（次スライド）
\end{frame}

\againframe{mcf-LP}

\begin{frame}[t]{相補性条件}
    \small
    $G_\varphi$のポテンシャル$x: V \to \R$に対し，$y: A \to \R$を次のように定める．
    \setlength{\abovedisplayskip}{3pt}
    \begin{align*}
        y(a) := \begin{cases}
            -c(a) + x(j) - x(i) & (\varphi(a) > 0) \text{\footnotesize ・・・$a^{-1}$に対するポテンシャル条件のスラック} \\
            0 & (\text{otherwise})
        \end{cases}
    \end{align*}

    \structure{Claim} $(x, y)$は双対実行可能解．
    \pause

    {\footnotesize \structure{(∵)} 

    \structure{$y \geq 0$:}
    $\varphi(a) > 0$ のとき，$a^{-1}$が$G_\varphi$に存在するので，$x$がポテンシャルより
    \setlength{\abovedisplayskip}{3pt}
    \setlength{\belowdisplayskip}{3pt}
    \[
    x(i) - x(j) \leq \ell_\varphi(a^{-1}) = -c(a).
    \]
    よって，$y(a) \geq 0$. 
    
    \bigskip
    \structure{$x(j) - x(i) - y(a) \leq c(a)$:}
    $\varphi(a) > 0$のときは，$y$の定義から自明に制約を満たす．$\varphi(a) = 0$のときは，$a$が$G_\varphi$に存在するので，ポテンシャル条件より
    \[ 
        x(j) - x(i) \leq \ell_\varphi(a) = c(a)
    \] 
    $y(a) = 0$より，制約を満たす．
    \qed}
\end{frame}
    
\begin{frame}[t]{相補性条件}
    \postit{\footnotesize
    \begin{minipage}{0.40\linewidth}
        \setlength{\abovedisplayskip}{0pt}
    \begin{align*}
        y(a) := \begin{cases}
            -c(a) + x(j) - x(i) & (\varphi(a) > 0) \\
            0 & (\text{otherwise})
        \end{cases}
    \end{align*}
    \end{minipage}}

    \small
    \structure{Claim} $\varphi$と$(x, y)$は相補性条件を満たす:
    \setlength{\abovedisplayskip}{3pt}
    \setlength{\belowdisplayskip}{3pt}
    \begin{align*}
        \varphi(a) (c(a) + x(i) - x(j) + y(a)) &= 0 \quad (a = ij \in A) \\
        (u(a) - \varphi(a)) y(a) &= 0 \quad (a \in A)
    \end{align*}
    すなわち，$\varphi$は最小費用流．
    \pause\vfill
    {\footnotesize \structure{(∵)} 
    相補性条件の1つ目は$y$の定義より自明なので，2つ目を示す．
    
    \smallskip
    背理法．$\varphi(a) < u(a)$かつ$y(a) > 0$と仮定する．このとき，$y$の定義から$\varphi(a) > 0$である．よって，$a$と$a^{-1}$の両方が$G_\varphi$に存在するので，ポテンシャル条件より次が成立:
    \setlength{\abovedisplayskip}{3pt}
    \setlength{\belowdisplayskip}{3pt}
    \begin{align*}
    x(j) - x(i) &\leq c(a) \\
    x(i) - x(j) &\leq - c(a).
    \end{align*}
    ∴$x(j) - x(i) = c(a)$. ところが，これを$y$の定義に代入すると，$y(a) = 0$となってしまい矛盾．
    \qed}
\end{frame}

% \begin{frame}{相補性条件}
%     \postit{\footnotesize $x(j) - x(i) - y(a) \leq c(a)$ \quad ($a=ij \in A$) \\ 
%             \footnotesize $y(a) \geq 0$ \quad ($a \in A$)}%
%     \small
%     \begin{block}{}
%         $b$流$\varphi$, 双対実行可能解$(x, y) \in \R^V \times \R_+^A$が最適$\iff$
%         \begin{itemize}
%             \item $x(j) - x(i) - y(a) = c(a)$ \quad ($a=ij \in A: \varphi(a) > 0$) \\
%             \onslide<2->{$\implies$ $x(i) - x(j) \leq -c(a)$ \dotfill $a^{-1} = ji \in \Aphilow$のポテンシャル条件}
%             \item $y(a) = 0$ \quad ($a = ij \in A: \varphi(a) < u(a)$) \\
%             \onslide<2->{$\implies$ $x(j) - x(i) \leq c(a)$ \dotfill $a = ij \in \Aphiup$のポテンシャル条件}
%             \item 
%             $\varphi(a) = 0$ \qquad ($a = ij \in A: x(j) - x(i) - y(a) < c(a)$) \\
%             $\varphi(a) = u(a)$ \quad ($a \in A: y(a) > 0$) \\
%             \onslide<2->{$\implies$ $a^{-1}, a$は残余ネットワークに存在しない}
%         \end{itemize}
%     \end{block}
    
%     \vfill
%     \onslide<3->
%     つまり$x$はポテンシャル．逆に，ポテンシャルから双対最適解$(x,y)$を構成できる．

%     \bigskip
%     ∴ $G_\varphi$にポテンシャルが存在 $\iff$ $\varphi$と相補性条件を満たす双対実行可能解$(x,y)$が存在 
%      $\iff$ $\varphi$は最適解 \qed
% \end{frame}

% \begin{frame}{流の分解定理}
%     $V^+ = \{i \in V : b(i) > 0\}$, $V^- = \{i \in V : b(i) < 0\}$
%     \begin{theorem}
%         $b$流$\varphi$は$V^+$--$V^-$パスと閉路の錐結合で表せる．すなわち，$\varphi = \sum_{k=1}^p \lambda_k \ones_{P_k} + \sum_{\ell=1}^q \mu_l \ones_{C_l}$ ($P_k$:$V^+$--$V^-$パス, $C_\ell$:閉路, $\lambda_k, \mu_\ell > 0$)
%     \end{theorem}
%     \structure{(証明)} まず循環流($b \equiv 0$)の場合を考える．
% \end{frame}

%%%%%%%%%%%%%%%%%%%%%%%%%%%%%%%%%%%%%%%%%%%%%
\section{アルゴリズム}
\againframe<3|handout:0>{outline}
%%%%%%%%%%%%%%%%%%%%%%%%%%%%%%%%%%%%%%%%%%%%%
\begin{frame}{アルゴリズム}
    以下のアルゴリズムを紹介する:

    \begin{itemize}
        \item 逐次最短路法
        \item 主双対法
        \item 容量スケーリング
    \end{itemize}
    {\footnotesize cf. 最小重み二部マッチングに対する逐次最短路法・主双対法（第3回）}

    \bigskip
    \structure{仮定}
    \begin{itemize}
        \item $b$, $u$は整数値
        \item $c$は非負 \\
        {\footnotesize 例えば，Hitchcock型輸送問題に帰着した後，費用$c' := c + M \ones$ ($M$:十分大) を考えればよい．}
    \end{itemize}
\end{frame}

\begin{frame}{逐次最短路法}
    \small
    \structure{記号} 流$\varphi$に対し，
    \begin{align*}
            U^+ &= \{ i \in V : b(i) > (\partial\varphi)(i) \} \quad \dots\dots\text{供給が余っている頂点}  \\
            U^- &= \{ i \in V : b(i) < (\partial\varphi)(i) \} \quad \dots\dots\text{需要が満たされていない頂点}     
    \end{align*}

    \begin{block}{逐次最短路法}
        \begin{algorithmic}[1]
            \STATE $\varphi \gets 0$ \COMMENT{初期流}
            \WHILE{$\partial\varphi \neq b$}
            \IF{残余ネットワーク$G_\varphi$上で$U^+$から$U^-$へ到達可能}
                \STATE 枝長$\ell_\varphi$に関する$G_\varphi$上の$U^+$--$U^-$最短路$P$を求める \\ \COMMENT{Bellman--Ford法で$O(mn)$時間}
                \STATE $\eps := \min\{\min_{a \in P} c_\varphi (a), b(s) - \partial\varphi(s), \partial\varphi(t) - b(t)\}$ \quad ({\footnotesize $s, t$: $P$の始点・終点})
                \STATE $P$に沿って$\varphi$を$\eps$だけ更新する．
            \ELSE
                \RETURN "$b$流は存在しない"
            \ENDIF
            \ENDWHILE
            \RETURN $\varphi$
        \end{algorithmic}
    \end{block}

\end{frame}

\begin{frame}[t]{逐次最短路法の証明}
    \small
    \begin{lemma}
        逐次最短路法で求まる流を$0 = \varphi_0, \varphi_1, \cdots,\varphi_\mu$とする．このとき，各$k=0, 1, \dots, \mu$に対し，$\varphi_k$は最小費用$b'$流である（{\footnotesize ここで，$b' = \partial \varphi_k$}）．
    \end{lemma}

\structure{(証明)} 
$k$に関する帰納法．$k=0$のときは，$c$が非負かつ$\varphi_0 = 0$より，自明に最小費用$0$流． 

\bigskip
$\varphi$を最小費用$b'$流とし，$G_\varphi$で$U^+$から$U^-$へ到達可能とする．

\bigskip
\onslide<+->
\structure{Claim.} $(G_\varphi, \ell_\varphi)$は$U^+$--$U^-$最短路をもつ \\
\only<+>{\footnotesize \structure{(∵)} 負閉路最適性条件より，$(G_\varphi, \ell_\varphi)$には負閉路が存在しない． \qed}

\bigskip
\onslide<+->
$P$を$U^+$--$U^-$最短路とし，$\varphi^+$を更新後の$b''$流とする．

\bigskip
\structure{Claim.} $\varphi^+$は最小費用$b''$流．\\
\only<+>{\footnotesize \structure{(∵)} $G_{\varphi^+}$に負閉路が存在したと仮定．更新前の$G_\varphi$には負閉路がなかったので，図のような状況．すると，迂回路$P'$は$P$より長さが短いことになり矛盾． \qed}
\end{frame}

\begin{frame}[t]{逐次最短路法の証明}
    \small
    \begin{lemma}
        逐次最短路法の途中で，$U^+, U^- \neq \emptyset$かつ$G_\varphi$で$U^+$から$U^-$へ到達不能になったとする．このとき，$(G, u)$に$b$流は存在しない．
    \end{lemma}

    \structure{(証明)} 
    $X \subseteq V$を$G_\varphi$において$U^+$から到達可能な頂点全体とする．定義から$U^+ \subseteq X \subseteq V \setminus U^-$.
    このとき，
    \begin{align*}
        b(X) 
        &= b(U^+) + b(X \setminus U^+) \\
        &> (\partial\varphi)(U^+) + (\partial\varphi)(X \setminus U^+) \\
        &= (\partial\varphi)(X) \\
        &= \varphi(\Out_G(X)) - \varphi(\In_G(X)) \\
        &= u(\Out_G(X)) - 0
    \end{align*}
    となり，Galeの定理より$b$流は存在しない．\qed
\end{frame}

\begin{frame}{逐次最短路法の計算量}

    \structure{計算量解析}
    \begin{itemize}
        \item 整数性より，$\varphi$の更新は高々$B$回 \\ 
        {\footnotesize \hfill ($B := \frac{1}{2} \sum_{i \in V}\abs{b(i)}$ ... 総供給/総需要量) }

        \item 各更新は1回の$(G_\varphi, \ell_\varphi)$上の最短路問題 \\
        \hfill $\longrightarrow$ Bellman--Ford法で$O(mn)$時間
    \end{itemize}
    {\footnotesize \hfill ($n = |V|, m = |A|$)}
    
    \begin{theorem}[伊理~(1960)]
        逐次最短路法は，$O(mnB)$時間で，最小費用$b$流を求めるか，$b$流が存在しないことを判定する．
    \end{theorem}
    \hfill \alert{擬多項式時間}
\end{frame}

\begin{frame}{主双対法}
    \small
    流$\varphi$と$(G_\varphi, \ell_\varphi)$のポテンシャル$p$の両方を保持することで，逐次最短路法の計算量を改善できる．

    \begin{definition}[簡約枝長 (復習)]
        枝長$\ell: A \to \R$のポテンシャル$p: V \to \R$に関する\alert{簡約枝長}$\ell_p: A \to \R$を次のように定める:
        \begin{align*}
       \ell_p(a) := \ell(a) + p(i) - p(j) \quad (a = ij \in A) 
        \end{align*}
    \end{definition}

    \bigskip
    \structure{観察}
    \begin{itemize}
        \item 
    任意の$s$--$t$路$P$に対し，
    \[
        \ell_p(P) = \ell(P) + p(s) - p(t).
    \]
    特に，$\ell_p$に関する$s$--$t$最短路 $\iff$ $\ell$に関する$s$--$t$最短路．
        \item $\ell_p \geq 0$ \dotfill\alert{Bellman-Ford法ではなくDijkstra法が使える！}
    \end{itemize}
\end{frame}

\begin{frame}{主双対法}
    
    \small
    \begin{block}{主双対法}
        \begin{algorithmic}[1]
            \STATE $\varphi \gets 0$ \COMMENT{初期流}
            \STATE $p \gets$ $(G, c)$のポテンシャル \COMMENT{初期ポテンシャル．Dijkstra法で$O(m + n \log n)$時間}
            \WHILE{$U^+, U^- \neq \emptyset$}
            \IF{残余ネットワーク$G_\varphi$で$U^+$から$U^-$へ到達可能}
                \STATE 簡約枝長$\ell_{\varphi,p}$に関する$G_\varphi$上の$U^+$--$U^-$最短路$P$と最短路長関数$d$を求める 
                \COMMENT{Dijkstra法で$O(m + n\log n)$時間}
                \STATE $\eps := \min\{\min_{a \in P} c_\varphi (a), b(s) - \partial\varphi(s), \partial\varphi(t) - b(t)\}$ \quad ({\footnotesize $s, t$: $P$の始点・終点})
                \STATE $P$に沿って$\varphi$を$\eps$だけ更新し，ポテンシャルを$p \gets p + d$と更新する．
            \ELSE
                \RETURN "$b$流は存在しない"
            \ENDIF
            \ENDWHILE
            \RETURN $\varphi$
        \end{algorithmic}
    \end{block}
\end{frame}

\begin{frame}{主双対法の証明}
    \small
    次の補題を示せば，あとは逐次最短路法の定理から従う．
    
    \begin{lemma}
    主双対法における$p$は，常に$(D_\varphi, \ell_\varphi)$上のポテンシャル．
    \end{lemma}
    \pause
    \structure{(証明)} 
    $p^+ := p + d$とする．{\footnotesize ($d: V \to \R$ ... $(G_\varphi, \ell_{\varphi,p})$上の（$U^+$からの）最短路長関数)}

    \bigskip
    $d$は$(G_\varphi, \ell_{\varphi,p})$のポテンシャルなので，任意の枝$a = ij \in A_\varphi$に対し，
    \begin{align*}
        d(j) - d(i) \leq \ell_{\varphi,p}(a) = \ell_\varphi(a) + p(i) - p(j)
    \end{align*}
    が成り立つ．移項すると $p^+(j) - p^+(i) \leq \ell_\varphi(a)$ となり，$p^+$は$(D_\varphi, \ell_\varphi)$のポテンシャルである．

    \pause
    \bigskip
    次に，$P$に沿って$\varphi$を更新した後も$p^+$がポテンシャルであることを示す．
    最短路$P$上の枝$a = ij$では，上の不等式で等号が成立:
    \begin{align*}
        d(j) - d(i) = \ell_{\varphi,p}(a) \quad \iff \quad p^+(j) - p^+(i) = \ell_\varphi(a)
    \end{align*}
    よって，更新後に生じる$P$の反転枝に関しても，ポテンシャル条件が成立． \qed
\end{frame}

\begin{frame}{主双対法の計算量}
    
    \structure{計算量解析}
    \begin{itemize}
        \item 整数性より，$\varphi$の更新は高々$B$回 \\ 
        {\footnotesize \hfill ($B := \frac{1}{2} \sum_{i \in V}\abs{b(i)}$ ... 総供給/総需要量) }

        \item 各更新は1回の$(G_\varphi, \ell_{\varphi, p})$上の最短路問題 \\
        \hfill $\longrightarrow$ Dijkstra法で$O(m + n \log n)$時間
    \end{itemize}
    {\footnotesize \hfill ($n = |V|, m = |A|$)}
    
    \vfill
    \begin{theorem}[Edmonds--Karp~(1970), 冨澤~(1971)]
        主双対法は，$O((m + n\log n)B)$時間で，最小費用$b$流を求めるか，$b$流が存在しないことを判定する．
    \end{theorem}
    \hfill \alert{擬多項式時間}
\end{frame}

\begin{frame}{容量スケーリング法}
    \small
    逐次最短路法・主双対法の反復回数を$O(B)$→$O(n \log b_{\max})$に改善したもの． \\
    {\hfill\footnotesize $B = \frac{1}{2} \sum_{i \in V}\abs{b(i)}$, $b_{\max} = \max_{i \in V} \abs{b(i)}$} \\
    \hfill \alert{弱多項式時間}

    \vfill
    \structure{追加の仮定} \\
    $u \equiv +\infty$ （{\footnotesize 例えば，Hitchcock型輸送問題に帰着すればよい}）

    \vfill
    \bigskip
    \structure{記号} 
    \begin{itemize}
        \item $\Delta \in \Z_{++}$: スケーリングパラメータ
        \item 流$\varphi$に対し，
        \begin{align*}
            U^+_\Delta &:= \{i \in U^+ : b(i) \geq (\partial\varphi)(i) + \Delta \} \subseteq U^+ \\
            U^-_\Delta &:= \{i \in U^- : b(i) \leq (\partial\varphi)(i) - \Delta \} \subseteq U^-
        \end{align*}
        \dotfill $\Delta$以上供給が余っている/需要が満たされていない頂点
    \end{itemize}
    整数性より$\Delta = 1$のときは，$(U^+_\Delta, U^-_\Delta) = (U^+, U^-)$.
\end{frame}

\begin{frame}{容量スケーリング法}
    \small
    \begin{block}{}
        \begin{algorithmic}[1]
            \STATE $\varphi \gets 0$ \COMMENT{初期流}
            \STATE $p \gets$ $(G, c)$のポテンシャル \COMMENT{初期ポテンシャル．Dijkstra法で$O(m + n \log n)$時間}
            \STATE $\Delta \gets 2^{\lceil \log_2 b_{\max} \rceil}$ \COMMENT{初期スケーリングパラメータ}
            \WHILE{$\Delta \geq 1$}
            \IF{残余ネットワーク$G_\varphi$で$U^+_\Delta$から$U^-_\Delta$へ到達可能}
                \STATE 簡約枝長$\ell_{\varphi,p}$に関する$G_\varphi$上の$U^+_\Delta$--$U^-_\Delta$最短路$P$と最短路長関数$d$を求める 
                \COMMENT{Dijkstra法で$O(m + n\log n)$時間}
                \STATE $P$に沿って$\varphi$を$\Delta$だけ更新し，ポテンシャルを$p \gets p + d$と更新する．
            \ELSE
                \STATE $\Delta \gets \Delta / 2$ \COMMENT{スケーリングパラメータの更新}
            \ENDIF
            \ENDWHILE
            \IF{$\partial \varphi = b$}
                \RETURN $\varphi$
            \ELSE
                \RETURN "$b$流は存在しない"
            \ENDIF
        \end{algorithmic}
    \end{block}
\end{frame}

\begin{frame}{容量スケーリング法の解析}
    $\Delta$が一定の間を\alert{$\Delta$フェーズ}と呼ぶ．

    \bigskip
    \begin{itemize}
        \setlength{\itemsep}{1em}

        \item $\Delta$フェーズにおける$\varphi(a)$は$\Delta$の倍数． \\
        {\footnotesize \structure{(∵)} 初期流は$\varphi \equiv 0$より自明．$\Delta$フェーズにおける更新では，毎回$\Delta$だけ増減するのでOK．最後に，$\Delta \rightarrow \Delta/2$フェーズに移行したときも，$\varphi$が$\Delta$の倍数なので，自明に$\Delta/2$の倍数でもあるのでOK．}

        \item $\varphi$は容量制約$0 \leq \varphi \leq u$を満たす． \\
        {\footnotesize \structure{(∵)} $u \equiv +\infty$より上限は自明．$\varphi(a) > 0$のとき，$\varphi(a)$は$\Delta$の倍数なので，更新で$\Delta$だけ減らしても負にはならない．}
    \end{itemize}
    


\end{frame}

\begin{frame}[t]{容量スケーリング法の解析}
    \small 
    \begin{lemma}
        各$\Delta$フェーズにおける更新は高々$n$回．
    \end{lemma}

    \structure{(証明)} $\varphi$を$\Delta$フェーズの開始時の流とする．
    $X$を$G_\varphi$において$U^+_{2\Delta}$から到達可能な頂点全体とすると，$U^+_{2\Delta} \subseteq X \subseteq V \setminus U^-_{2\Delta}$. また，$u \equiv +\infty$より，$G$に$X$から出る枝は存在しない．

    \pause
    \bigskip
    $(s_1, t_1), \dots, (s_k, t_k)$: $\Delta$フェーズ中の更新で使われた$P$の始点・終点（重複を許す）
    
    \bigskip
    \structure{Claim} $\abs{\{\ell : s_\ell \in X, t_\ell \notin X\}} \leq \abs{\{ \ell : s_\ell \notin X, t_\ell \in X\}}$ \\
    \pause
    {\footnotesize 
    \structure{(∵)} $G$に$X$の外に出る枝はないので，$X$から出るパスを使うには，それより前に$X$へ入るパスを使い，$G_\varphi$に$X$から出る枝を逆向き枝として作る必要がある． \qed}
\end{frame}

\begin{frame}[t]{容量スケーリング法の解析}
    \small
    \structure{Claim} $\abs{\{\ell : s_\ell \in X, t_\ell \notin X\}} \leq \abs{\{ \ell : s_\ell \notin X, t_\ell \in X\}}$ 

    \bigskip
    よって，
    \begin{align*}
        k
        &\leq \abs{\{\ell : s_\ell, t_\ell \in X\}} + 2\abs{\{\ell : s_\ell \notin X, t_\ell \in X\}} + \abs{\{\ell : s_\ell, t_\ell \notin X\}} \\
        &= \abs{\{\ell: t_\ell \in X\}} + \abs{\{\ell: s_\ell \notin X\}} \\
        &\leq \abs{X \cap U^-_\Delta} + \abs{U^+_{\Delta} \setminus X} \\
        &\leq n. \qed
    \end{align*}

    \pause
    \begin{theorem}[Edmonds--Karp~(1970)]
        容量スケーリング法は，$O(n(m + n \log n) \log b_{\max})$時間で，最小費用$b$流を求めるか，$b$流が存在しないことを判定する．
    \end{theorem}
    {\footnotesize \hfill ($n = |V|, m = |A|, b_{\max} = \max_{i \in V} \abs{b(i)}$)} \\
    \hfill \alert{弱多項式時間}
\end{frame}


\begin{frame}{最小費用流: まとめ}
    \small
    
    \structure{まとめ}
    \begin{itemize}
        \item 最小費用流問題の定義，Galeの定理，負閉路最適性条件
        \item 逐次最短路法，主双対法，容量スケーリング法
    \end{itemize}

    \bigskip
    \structure{その他のアルゴリズム}
    \begin{itemize}
        \item 最小平均閉路消去法 $O(m^3 n^2 \log n)$時間 \alert{強多項式時間} (Goldberg--Tarjan~(\emph{JACM} 1989))
        \item 費用スケーリング $O(mn \log(n^2/m) \log (nC))$時間  {\hfill\footnotesize $C = \max_{a \in A} \abs{c(a)}$} \\
        (Goldberg--Tarjan~(\emph{Math of OR} 1990))
        \item 現在最速 $m^{1 + o(1)}$時間 (\emph{JACM} 2025\footnote{L. Chen, R. Kyng, Y. P. Liu, R. Peng, M. P. Gutenberg, and S. Sachdeva, \textit{``Maximum Flow and Minimum-Cost Flow in Almost-Linear Time''}, \textit{JACM}, 2025.})
    \end{itemize}
\end{frame}
%%%%%%%%%%%%%%%%%%%%%%%%%%%%%%%%%%%%%%%%%%%%%
\end{document}